\clearpage
\VerdierPhysicalPageAnchor{0144}%
\begingroup
\small
\setlength{\parskip}{0pt}%
\setlength{\emergencystretch}{1.5em}%
\setlength{\abovedisplayskip}{3pt plus 1pt minus 1pt}%
\setlength{\belowdisplayskip}{3pt plus 1pt minus 1pt}%

\VerdierSemanticAnchor{ch2-2-3-1-proof-c-bis-converse}
\noindent 飽和である。d) の証明により、部分圏
$Q(Q^{-1}(\mathcal C))$ は厳密充満かつ飽和である。一方、圏
$Q(Q^{-1}(\mathcal C))$ は $\mathcal C$ の厳密充満閉包である。したがって
$\mathcal C$ は飽和である。これで主張 c)$^{\mathrm{bis}}$ が示された。

\VerdierSemanticAnchor{ch2-2-3-1-proof-d-bis}
最後の主張を示す。$\mathcal C$ が厳密充満ならば
$Q(Q^{-1}(\mathcal C))=\mathcal C$ であることは明らかである。
$\widehat{\mathcal B}$ を含む厳密充満三角圏 $\mathcal A$ をとる。
$\mathcal A=Q^{-1}(Q(\mathcal A))$ を示せばよい。$\mathcal B$ を
$\widehat{\mathcal B}$ で置き換えても圏 $\mathcal D/\mathcal B$ は変わらない
(2.2.11) ので、$\mathcal B$ は飽和で
$\mathcal B=\widehat{\mathcal B}$ と仮定してよい。包含
$\mathcal A\subset Q^{-1}(Q(\mathcal A))$ は明らかなので、
$\mathcal D/\mathcal B$ で $\mathcal A$ の対象と同型になる
$\mathcal D$ の任意の対象が $\mathcal A$ の対象であることを示せばよい。

\VerdierSemanticAnchor{ch2-2-3-1-proof-d-bis-roof-setup}
$X$ を $\mathcal A$ の対象、$Y$ を $\mathcal D$ の対象とし、同型
$m\colon Q(X)\xrightarrow{\sim}Q(Y)$ をとる。射 $m$ は
$m=Q(f)Q(s)^{-1}$ と分解でき、ここで
$s\in S_{\mathcal D}(\mathcal B)$ (2.2.2)、$Q(f)$ は同型である。
$\mathcal B$ は飽和なので $f\in S_{\mathcal D}(\mathcal B)$ である
（(2.1.8) と (2.2.4)）。したがって $\mathcal D$ に次の図式がある：

\VerdierSemanticAnchor{ch2-2-3-1-proof-d-bis-roof}
\begin{center}
  \begin{tikzcd}[column sep=4.2em,row sep=2.5em,baseline=(current bounding box.center)]
    & Z \arrow[dl,"s"'] \arrow[dr,"f"] & \\
    X && Y
  \end{tikzcd}
  \qquad $s,f\in S_{\mathcal D}(\mathcal B)$\quad .
\end{center}

\VerdierSemanticAnchor{ch2-2-3-1-proof-d-bis-complete}
$s$ は $S_{\mathcal D}(\mathcal B)$ に属するので、$W$ が $\mathcal B$ の対象である
完全三角 $(Z,X,W,s,u,v)$ が存在する (2.1.7)。したがって $W$ は
$\mathcal A$ の対象であり、$\mathcal A$ は厳密充満だから $Z$ も
$\mathcal A$ の対象である (1.2.4)。同じ議論により $Y$ も
$\mathcal A$ の対象となり、命題の証明が完了する。

\VerdierSemanticAnchor{ch2-2-3-2}
\medskip\noindent\textbf{2.3.2.}\quad
\VerdierSemanticAnchor{ch2-2-3-2-free-object}
$F\colon C\longrightarrow C'$ を関手、$X$ を $C'$ の対象とする。
$u\colon F(Y)\longrightarrow X$ が存在し、任意の射
$F(Z)\xrightarrow{v}X$ に対して下図を可換にする唯一の射
$w\colon Z\longrightarrow Y$ が存在するとき、$C$ の対象 $Y$ は
\emph{$X$ 上 $F$-右自由}、または単に\emph{$X$ 上右自由}であるという：

\VerdierSemanticAnchor{ch2-2-3-2-free-object-diagram}
\begin{center}
  \begin{tikzcd}[column sep=4.7em,row sep=2.7em]
    F(Y) \arrow[rr,"u"] && X \\
    & F(Z) \arrow[ul,"{F(w)}"] \arrow[ur,"v"'] &
  \end{tikzcd}
  \quad .
\end{center}

\VerdierSemanticAnchor{ch2-2-3-2-freedom-arrow-boundary}
矢 $u\colon F(Y)\longrightarrow X$ を、組 $(X,Y)$ の（関手 $F$ に関する）
\emph{自由矢}という。これはまた、ある
\endgroup
